\scanpage{37}
% L1–2
\deleted{Enfin on conclut que tt morphisme projectif [simple] est $\in M$ \unsure{grâce} à \unsure{e)}.}

% L3–4
\deleted{\uline{Corollaire} Supposons \deleted{de plus} \inserted{au lieu de c)} la condition suivante satisfaite :}

% L5–11
\deleted{c') Soit \unsure{$g$|$\varnothing$} : $X'\xrightarrow{f}X\xrightarrow{g}Y$ \inserted{des} morphismes \deleted{\unsure{de préschémas}} propres, $X_0$ un sous-préschéma fermé de $X$, $U=X-X_0$, $U'=f^{-1}(U)$. \deleted{et} Supposons que $U'\unsure{\mathord{-}|\to}U$ soit un isomorphisme, \deleted{et} que $X_0\xrightarrow{g_0}Y$ \deleted{soit} et $gf:X'-Y$ soient $\in M$. Alors $g\in M$.}

% L12–13
\deleted{Alors tt morphisme propre $f:X-Y$ \deleted{est} $\in M$.}

% L14–15
\deleted{En effet, il est trivial que c') implique \unsure{a)|c)}, donc tt morphisme projectif est $\in M$.}

% L16–22
\deleted{On obtient le cas d'un morphisme \deleted{\unsure{quelcon|quelconque}} propre quelconque en utilisant \unsure{e)|c)}, \unsure{[}qui permet \deleted{\unsure{de|du}} supposer $Y$ affine,\unsure{]} \deleted{et c'} et d'utiliser \inserted{la récurrence} \deleted{le lemme de} \deleted{Chow, puis c'} noethérienne. Le lemme de Chow donne \deleted{\unsure{donc}} joint à c') donne alors le résultat voulu.}

% archive sheet number: 36
